\documentclass[12pt,a4paper,oneside]{article}
\usepackage[brazil]{babel}
\usepackage[utf8]{inputenc}
\usepackage[dvipsnames]{xcolor}
\usepackage{listings}
\usepackage{wrapfig}
\usepackage{olymp}

\setcounter{problem}{1}

\begin{document}

\begin{problem}{IMC}{imc}{2}

\begingroup
\setlength{\intextsep}{0pt}%
% \setlength{\columnsep}{0pt}


Escreva um programa que calcula o IMC (Índice de Massa Corporal) de uma pessoa.

O IMC é calculado da seguinte forma:

$$\text{IMC} = \frac{\text{massa}}{\text{altura}^2}$$

sendo a massa dada em quilogramas e a altura em metros.

Por exemplo, uma pessoa com $1.80$m de altura e $75$Kg tem um IMC de valor $75/(1.80)^2 = 23.15$.

% \Notes

\InputFile

A entrada contém apenas uma linha, com dois valores inteiros, a altura $A$ em cm e a massa $M$ em Kg.

Restrições: $100 \leq A \leq 220$, $30 \leq M \leq 200$.

\endgroup

\OutputFile

Escreva o valor do IMC com 2 casas decimais.

% \Notes

\Example

\begin{example}
\exmp{
180 75
}{
23.15
}%
\end{example}

\begin{example}
\exmp{
180 80
}{
24.69
}%
\end{example}

\begin{example}
\exmp{
160 100
}{
39.06
}%
\end{example}

%Comentários sobre o exemplo, se necessário.

\end{problem}

\end{document}
